\section{Diseño del Sistema}
\label{sec:diseno-sistema}

Esta sección describe el diseño del prototipo AeroTrack, la consola con la que se configura y
opera el sistema. El apartado~\ref{subsec:arquitectura} presenta los seis componentes que
organizan el procesamiento de las imágenes y el motor de identidad entre cámaras. Los
apartados siguientes documentan la consola: la gestión de usuarios por roles, la gestión de
proyectos independientes, el asistente de configuración y las alertas de seguridad con aviso
por correo.

\subsection{Arquitectura propuesta}
\label{subsec:arquitectura}

El prototipo organiza el procesamiento en seis componentes, cada uno con una
responsabilidad única (Figura~\ref{fig:arquitectura}). Todo se ejecuta en el equipo local: el
servidor solo acepta conexiones del propio equipo y ningún fotograma sale de él.

\begin{figure}[H]
\centering
\resizebox{\anchofigura}{!}{%
\begin{tikzpicture}[
  comp/.style={draw=black,thick,rounded corners=2pt,align=center,font=\small,
               minimum width=7.6cm,minimum height=1.25cm,inner sep=4pt},
  dato/.style={draw=black,dashed,rounded corners=2pt,align=center,font=\small,
               minimum width=4.6cm,minimum height=1.1cm,inner sep=4pt},
  flecha/.style={-{Stealth[length=2.2mm]},thick}]
  \node[comp] (a) at (0,0) {\textbf{1. Adquisición}\\archivo de video, cámara USB o flujo de red, por cámara};
  \node[comp,below=0.55cm of a] (d) {\textbf{2. Detección}\\YOLO11 (variante según hardware) y filtros de\\zona útil, tamaño relativo e inmovilidad};
  \node[comp,below=0.55cm of d] (s) {\textbf{3. Seguimiento por cámara}\\ByteTrack en dos etapas sobre el punto de los pies,\\con caja predicha y firma de color del torso};
  \node[comp,below=0.55cm of s] (c) {\textbf{4. Calibración y proyección}\\homografía por cámara hacia el plano (opcional)};
  \node[comp,below=0.55cm of c] (i) {\textbf{5. Identidad entre cámaras}\\vector OSNet de cuerpo entero, tramos, compuerta de\\tiempo y plano, reagrupación al cerrar la sesión};
  \node[comp,below=0.55cm of i] (p) {\textbf{6. Persistencia, analítica y presentación}\\conteos, alertas, reportes y consola web};
  \node[dato,right=1.0cm of i] (m) {Memoria de apariencia\\(base local, 24 h por defecto)};
  \node[dato,right=1.0cm of p] (g) {Grabación de la sesión\\(identificadores y posiciones)};
  \foreach \x/\y in {a/d,d/s,s/c,c/i,i/p} { \draw[flecha] (\x) -- (\y); }
  \draw[flecha,{Stealth[length=2.2mm]}-{Stealth[length=2.2mm]}] (i) -- (m);
  \draw[flecha] (p) -- (g);
\end{tikzpicture}}
\caption{Arquitectura del prototipo AeroTrack. Los recuadros discontinuos son los datos por
persona que el sistema conserva en el equipo.}
\label{fig:arquitectura}
\end{figure}

El primer componente es la adquisición, que lee cada fuente de cámara, sea un archivo de
video, una cámara USB o un flujo de red, y la abre una sola vez. El segundo es la detección: el
detector YOLO11 (Jocher y Qiu, 2024) localiza personas en cada muestra con una caja y una
confianza. La variante del detector no la elige el operador, sino el sistema según el hardware
del equipo: nano en CPU, y pequeña o mediana con GPU. Tras la detección se aplican tres filtros
contra falsos positivos: se descarta lo que cae fuera de la zona útil de la imagen o del área de
trabajo del plano, las detecciones dudosas mucho más pequeñas que las personas seguras de la
misma cámara (gorros, carteles) y los objetos que nunca se mueven y se detectan con baja
confianza.

El tercero es el seguimiento dentro de cada cámara. Aplica la asociación en dos etapas de
ByteTrack (Zhang et al., 2022) sobre el punto de los pies, que es el que la homografía puede
proyectar, con un filtro de Kalman de velocidad constante: primero se emparejan las
detecciones de alta confianza y después las de baja confianza, para no perder a una persona
parcialmente ocluida. Dentro de la puerta espacial, las coincidencias se ordenan además con
dos señales al estilo de BoT-SORT (Aharon et al., 2022): la superposición entre la caja predicha
y la detectada, y una firma de color de la ropa calculada sobre dos franjas del torso.

El cuarto componente es la calibración y proyección: cada cámara tiene una homografía propia
hacia el plano, calculada con al menos cuatro correspondencias entre el suelo del video y el
plano, que se valida por consistencia geométrica antes de aceptarse. La calibración es
opcional; sin ella la cámara cuenta y sigue personas, pero no las ubica en el plano.

El quinto es el motor de identidad entre cámaras (Figura~\ref{fig:identidad}). Para cada vista
confiable de una persona calcula un vector de apariencia OSNet (Zhou et al., 2019) de 512
dimensiones sobre el recorte de cuerpo entero, con las partes tapadas por otra persona
rellenadas de gris, y lo acumula en el tramo continuo de seguimiento al que pertenece. Una
persona lleva un identificador provisional hasta reunir suficientes vistas y solo entonces recibe
un identificador confirmado, que es el único que cuenta en los conteos. La unión entre cámaras
exige a la vez parecido de apariencia y coherencia de tiempo y, si hay calibración, de posición
en el plano, y solo se intenta entre cámaras que el operador declaró relacionadas y cuyos relojes
declaró sincronizados; ante una candidata ambigua, la asociación se marca como incierta en
lugar de forzarse. Con la memoria de apariencia activa, que es la configuración por defecto, el
motor compara además cada identidad nueva con todas las personas guardadas en una base
de datos local del proyecto durante el plazo de retención (24 horas por defecto, hasta 168), de
modo que reconoce a la misma persona en sesiones y días posteriores. Al cerrar la sesión,
reagrupa los tramos con la grabación completa a la vista y renumera a las personas de 1 a N.

\begin{figure}[H]
\centering
\resizebox{\anchofigura}{!}{%
\begin{tikzpicture}[
  paso/.style={draw=black,thick,rounded corners=2pt,align=center,font=\small,
               minimum width=6.4cm,minimum height=1.0cm,inner sep=3pt},
  decision/.style={draw=black,thick,diamond,aspect=2.6,align=center,font=\small,inner sep=1pt},
  fin/.style={draw=black,thick,rounded corners=6pt,align=center,font=\small,
              minimum width=3.6cm,minimum height=0.9cm,inner sep=3pt},
  flecha/.style={-{Stealth[length=2.2mm]},thick}]
  \node[paso] (n1) at (0,0) {Vista confiable de una persona seguida\\(confianza, altura, borde, oclusión)};
  \node[paso,below=0.5cm of n1] (n2) {Vector OSNet de cuerpo entero (512 dimensiones)\\y prototipo del tramo de seguimiento};
  \node[decision,below=0.5cm of n2] (d1) {¿Vistas suficientes\\para confirmar?};
  \node[fin,left=1.2cm of d1] (f0) {Identificador\\provisional};
  \node[decision,below=0.6cm of d1] (d2) {¿Coincide con una\\persona guardada?};
  \node[fin,right=1.2cm of d2] (f1) {Recupera su ID\\(reidentificada)};
  \node[decision,below=0.6cm of d2] (d3) {¿Candidata en cámara\\relacionada, coherente en\\apariencia, tiempo y plano?};
  \node[fin,left=1.2cm of d3] (f3) {ID nuevo\\(local)};
  \node[decision,below=0.6cm of d3] (d4) {¿Es inequívoca?};
  \node[fin,right=1.2cm of d4] (f2) {Reutiliza su ID\\(estimada)};
  \node[fin,left=1.2cm of d4] (f4) {No fuerza la unión\\(incierta)};
  \draw[flecha] (n1) -- (n2);
  \draw[flecha] (n2) -- (d1);
  \draw[flecha] (d1) -- node[above,font=\small]{no} (f0);
  \draw[flecha] (d1) -- node[right,font=\small]{sí} (d2);
  \draw[flecha] (d2) -- node[above,font=\small]{sí} (f1);
  \draw[flecha] (d2) -- node[right,font=\small]{no} (d3);
  \draw[flecha] (d3) -- node[above,font=\small]{no} (f3);
  \draw[flecha] (d3) -- node[right,font=\small]{sí} (d4);
  \draw[flecha] (d4) -- node[above,font=\small]{sí} (f2);
  \draw[flecha] (d4) -- node[above,font=\small]{no} (f4);
  \node[font=\small\itshape,align=center,below=0.5cm of d4]
    {La consulta a personas guardadas solo ocurre con la memoria de apariencia activa,\\que es la configuración por defecto (RNF-01 no conforme).};
\end{tikzpicture}}
\caption{Decisión de identidad de una persona en el motor de identidad entre cámaras.}
\label{fig:identidad}
\end{figure}

El sexto componente reúne persistencia, analítica y presentación. Calcula la ocupación por
zona, el flujo por los accesos y las alertas de aglomeración a partir de las personas
confirmadas, y lo expone en la consola, que ofrece configuración, plano en vivo, video por
cámara, alertas y reportes. No se limita a mantener un estado agregado en memoria: guarda en
el equipo una grabación de cada sesión con, por muestra y por persona, el identificador, la caja,
la posición en el plano y su trayectoria reciente, sin plazo de retención automático, y a partir de
esas grabaciones calcula las métricas comerciales por negocio. No guarda fotogramas de las
cámaras en vivo ni vectores de apariencia en la grabación. Las consecuencias de esa
persistencia se examinan en el Anexo~\ref{anx:marco}, apartado B.9.

\subsection{Gestión de usuarios y roles}
\label{subsec:diseno-usuarios}

El acceso a la consola requiere usuario y contraseña. En la primera ejecución el sistema pide
crear un primer usuario (Figura~\ref{fig:primer-usuario}), que queda con el rol de operador y
puede dar de alta a los demás. La contraseña debe tener al menos seis caracteres y se guarda
en el equipo como resumen criptográfico con sal propia, nunca en texto plano.

\figura{aerotrack-01-primer-usuario.png}
  {Creación del primer usuario en la primera ejecución de la consola.}
  {fig:primer-usuario}

Una vez dentro, la vista general muestra en la cabecera el proyecto activo junto al usuario y su
rol, y organiza la consola en pestañas: vista general, alertas, videos y resultados, reportes y
configuración (Figura~\ref{fig:vista-general}). Desde esta vista se filtra por nivel, cámara, zona
de imagen y rendimiento, y se inicia el monitoreo sobre el plano del nivel. En la captura los
contadores de cámaras seleccionadas, entradas, salidas y tiempo de análisis están en cero
porque el proyecto aún no tiene cámaras configuradas.

\figura{aerotrack-02-vista-general.png}
  {Vista general de la consola, con el proyecto activo en la cabecera y el plano del nivel 3.}
  {fig:vista-general}

El sistema distingue dos roles (Figura~\ref{fig:usuarios}). El operador configura cámaras y
plano; el administrador consulta resultados y ajusta los umbrales de alerta, sin acceso a la
configuración. Los roles se aplican en el servidor y no solo en la interfaz. El panel de usuarios
lista las cuentas existentes con su rol, permite eliminarlas y permite crear otras indicando
nombre, contraseña y rol.

\figura{aerotrack-03-usuarios.png}
  {Panel de usuarios del sistema y formulario de alta con selección de rol.}
  {fig:usuarios}

\subsection{Gestión de proyectos}
\label{subsec:diseno-proyectos}

Un proyecto representa un espacio monitoreado, por ejemplo un nivel de un terminal. Cada
proyecto guarda por separado su propio plano, sus cámaras, su calibración, sus incidentes y su
memoria de apariencia, de modo que la configuración y los datos de un espacio no se mezclan
con los de otro. Al entrar en la configuración, la consola ofrece dos acciones: crear un proyecto
nuevo, que conduce al asistente descrito en el apartado~\ref{subsec:diseno-asistente}, o añadir
una cámara al proyecto abierto (Figura~\ref{fig:crear-proyecto}). Para crear el proyecto basta
con darle un nombre.

\figura{aerotrack-04-crear-proyecto.png}
  {Opciones para crear un proyecto nuevo o añadir una cámara al proyecto abierto.}
  {fig:crear-proyecto}

Los proyectos guardados se listan con su nombre, el espacio al que corresponden, el número
de cámaras y de planos, el momento de la última edición y su estado
(Figura~\ref{fig:proyectos-guardados}). En la captura figuran dos proyectos, ambos marcados
como incompletos: el proyecto principal, sobre el nivel 3 del aeropuerto y con cuatro planos, y
el proyecto Terminal~A, Nivel~1, con un plano. Desde la lista se abre un proyecto, se retoma su
configuración en el punto en que quedó, se renombra directamente sobre su tarjeta
(Figura~\ref{fig:renombrar}) o se elimina.

\figura{aerotrack-05-proyectos-guardados.png}
  {Lista de proyectos guardados, cada uno con sus cámaras, planos y estado.}
  {fig:proyectos-guardados}

\figura{aerotrack-06-renombrar-proyecto.png}
  {Cambio de nombre de un proyecto sobre su propia tarjeta.}
  {fig:renombrar}

Al entrar en la consola, una pantalla pregunta en qué proyecto se quiere trabajar, muestra para
cada uno su espacio, su número de cámaras y el estado de su configuración, y señala el último
abierto (Figura~\ref{fig:selector-proyecto}). Así el usuario elige de forma expresa el espacio en
el que va a trabajar y no mezcla datos de dos lugares distintos.

\figura{aerotrack-07-selector-proyecto.png}
  {Selección del proyecto de trabajo al entrar en la consola.}
  {fig:selector-proyecto}

\subsection{Asistente de configuración}
\label{subsec:diseno-asistente}

La configuración de un proyecto se completa con un asistente. Las capturas de este apartado
corresponden a su versión de ocho pasos: proyecto, plano, conexión, prueba, ubicación,
calibración, zonas y alertas, y revisión. Una barra superior muestra los pasos, marca los
completados y lleva la cuenta de cuántos están listos, y una indicación en la cabecera confirma
que el avance queda guardado, de modo que un proyecto incompleto puede retomarse desde la
lista de proyectos. Las capturas de los pasos 4, 6 y 8 se tomaron sin video, porque no se
dispone de grabaciones del terminal.

La versión vigente del asistente conserva las mismas tareas, pero las reúne en cinco etapas. La
Tabla~\ref{tab:asistente} muestra la correspondencia, y los pasos se describen a continuación
en el orden de las capturas.

\begin{table}[H]
\centering
\small
\caption{Correspondencia entre los pasos de las capturas y las etapas del asistente vigente.}
\label{tab:asistente}
\begin{tabular}{@{}L{5.2cm}L{9.4cm}@{}}
\toprule
\textbf{Etapa vigente} & \textbf{Pasos de las capturas que reúne} \\
\midrule
Proyecto & Paso 1 (datos del espacio y tipo de fuente) y paso 2 (plano) \\
Cámaras & Paso 3 (cámaras, fuente, zona útil y accesos) y paso 4 (prueba y vista en vivo) \\
Homografía y cámaras relacionadas & Paso 5 (ubicación) y paso 6 (calibración, ahora opcional), más la relación entre cámaras y la sincronización de sus relojes \\
Contexto comercial & Paso 7 (zonas, negocios y parámetros de aviso de aglomeración) \\
Validación & Paso 8 (revisión) \\
\bottomrule
\end{tabular}
\end{table}
\verificar{ubicación de la cámara sobre el plano (paso 5 de las capturas) dentro de la etapa de
homografía del asistente vigente: inferida de la estructura de la consola, sin captura nueva que
la muestre.}

\subsubsection*{Paso 1: proyecto}

Se identifica el espacio con el nombre del aeropuerto o proyecto y el del terminal y nivel, y se
elige el tipo de fuente: videos grabados, para analizar grabaciones y consultar resultados por
periodo, o cámaras en vivo, conectadas por USB, RTSP o IP (Figura~\ref{fig:paso1}). Una nota
lateral recuerda que el análisis es de presencia y recorridos, sin reconocimiento facial.

\figura{aerotrack-08-paso1-proyecto.png}
  {Paso 1 del asistente: datos del espacio y tipo de fuente.}
  {fig:paso1}

\subsubsection*{Paso 2: plano}

Se define el plano del espacio: se elige uno de los planos del aeropuerto, se importa una
imagen o un PDF, o se dibuja el espacio con rectángulos, polígonos y líneas. El plano se
acompaña de su ancho y su alto en metros, y admite marcar un área de trabajo como
referencia opcional (Figura~\ref{fig:paso2}).

\figura{aerotrack-09-paso2-plano.png}
  {Paso 2 del asistente: plano del nivel y sus dimensiones en metros.}
  {fig:paso2}

\subsubsection*{Paso 3: conexión}

Se añaden las cámaras del proyecto (Figura~\ref{fig:paso3}). Para cada una se indica el plano
asociado, el nombre visible en el mapa, la ubicación, la orientación y la fuente, con la ruta del
video o la dirección de la cámara, y se fijan el umbral de alerta en personas y el tiempo de
concentración sostenida (Figura~\ref{fig:paso3-camara}). En el mismo paso se marca la zona
útil de la imagen, que deja fuera reflejos, otras plantas y áreas que no deben analizarse, y, si la
cámara ve una puerta, un acceso formado por un área exterior y otra interior, entre las que se
cuenta el paso de personas (Figura~\ref{fig:paso3-zona}).

\figura{aerotrack-10-paso3-conexion.png}
  {Paso 3 del asistente: lista de cámaras del proyecto, todavía vacía.}
  {fig:paso3}

\figura{aerotrack-11-paso3-camara.png}
  {Paso 3 del asistente: datos de la cámara, fuente y umbrales de alerta.}
  {fig:paso3-camara}

\figura{aerotrack-12-paso3-zona-util.png}
  {Paso 3 del asistente: zona útil de la imagen y acceso de un local.}
  {fig:paso3-zona}

\subsubsection*{Paso 4: prueba}

Se comprueba la cámara en movimiento para revisar encuadre, luz y retardo. El panel muestra
el estado de la fuente, su resolución, su frecuencia de cuadros y si está calibrada, y permite
iniciar una prueba de seguimiento en modo laboratorio que presenta el video procesado con
identificadores temporales (Figura~\ref{fig:paso4}).

\figura{aerotrack-13-paso4-prueba.png}
  {Paso 4 del asistente: prueba de la cámara, con el visor sin señal.}
  {fig:paso4}

\subsubsection*{Paso 5: ubicación}

Se coloca la cámara sobre el plano en el punto donde está instalada y se orienta hacia donde
apunta, con su dirección, alcance, altura y ancho de cobertura (Figura~\ref{fig:paso5}). Esa
cobertura es orientativa: la precisión la aporta la calibración del paso siguiente.

\figura{aerotrack-14-paso5-ubicacion.png}
  {Paso 5 del asistente: posición y cobertura de la cámara sobre el plano.}
  {fig:paso5}

\subsubsection*{Paso 6: calibración}

Se relaciona el suelo del video con el plano marcando puntos de referencia, como esquinas o
cruces de baldosas, primero en la imagen y después en el mapa (Figura~\ref{fig:paso6}). Se
requieren al menos cuatro puntos repartidos; con cinco o más, cada referencia se comprueba
además dejándola fuera del ajuste. En la versión vigente la calibración es opcional, y la misma
etapa recoge qué cámaras están relacionadas y si sus relojes están sincronizados, condición sin
la cual el sistema no asocia personas entre cámaras.

\figura{aerotrack-15-paso6-calibracion.png}
  {Paso 6 del asistente: calibración del suelo del video con el plano.}
  {fig:paso6}

\subsubsection*{Paso 7: zonas y alertas}

Se dibujan sobre el plano las zonas de interés con polígonos, rectángulos o círculos, y se define
cuándo avisar de una aglomeración mediante un radio en metros, un número mínimo de
personas y un tiempo de permanencia (Figura~\ref{fig:paso7}). La pantalla advierte que una fila y
una aglomeración pueden parecerse geométricamente, por lo que los límites deben ajustarse al
uso de cada zona. En la versión vigente esta etapa recoge además los negocios del piso, que
alimentan el panel comercial.

\figura{aerotrack-16-paso7-zonas.png}
  {Paso 7 del asistente: trazado de una zona y parámetros de aviso de aglomeración.}
  {fig:paso7}

\subsubsection*{Paso 8: revisión}

Se revisa el proyecto antes de abrir el monitoreo: una lista indica qué está preparado y qué
sigue pendiente, y un resumen muestra el número de cámaras, zonas dibujadas y fuentes
probadas, junto con la escala (Figura~\ref{fig:paso8}). En la captura, con la cámara sin probar ni
calibrar, la opción de guardar y abrir el monitoreo aparece desactivada y queda disponible la
de continuar las pruebas en laboratorio. La pantalla aclara que una verificación completada no
sustituye la validación de precisión.

\figura{aerotrack-17-paso8-revision.png}
  {Paso 8 del asistente: revisión del proyecto antes de abrir el monitoreo.}
  {fig:paso8}

\subsection{Alertas de seguridad con aviso por correo}
\label{subsec:diseno-alertas}

La vista de alertas reúne lo que necesita el personal de seguridad para atender una
aglomeración (Figura~\ref{fig:alertas-correo}). El apartado \enquote{Cuándo avisar} fija los
umbrales que definen una aglomeración: el número mínimo de personas, la permanencia en
segundos y el radio en metros; en la captura, 4 personas, 3 segundos y 1,5 metros. Estos
umbrales se cambian desde la misma vista, sin entrar en la configuración, y son el único ajuste
que también puede hacer el administrador.

El panel \enquote{Avisar por correo}, desplegado en la captura y todavía desactivado, pide la
cuenta que envía, su contraseña de aplicación y los correos que reciben el aviso, separados por
coma, hasta un máximo de veinte, y permite enviar un correo de prueba. Una nota aclara que el
servicio de correo no acepta la contraseña normal de la cuenta, sino una contraseña de
aplicación generada con la verificación en dos pasos. Debajo, tres contadores muestran las
alertas activas, el total de la sesión y los incidentes sin revisar, y la bitácora de incidentes
conserva cada caso aunque la sesión termine, para que el personal lo marque como revisado,
falsa alarma o resuelto. En la captura la bitácora aparece sin registros porque no se han
procesado grabaciones del terminal.

Con el aviso activado, el sistema envía un correo por cada incidente nuevo de aglomeración,
con el lugar, la cámara y el motivo del aviso, y la advertencia de que requiere comprobación
humana; un aviso ya enviado no se reenvía y una misma zona no genera más de un correo cada
120 segundos. La contraseña se guarda solo en el equipo, fuera de los archivos del proyecto: la
interfaz nunca la muestra de vuelta y no se copia al duplicar un proyecto.

\figura{aerotrack-18-alertas-correo.png}
  {Vista de alertas: umbrales de aglomeración, aviso por correo desplegado, contadores y
   bitácora de incidentes.}
  {fig:alertas-correo}
